\begin{abstract}

Single-prompt code generation with Large Language Models (LLMs) frequently fails on tasks requiring multi-step reasoning or edge-case handling. Multi-agent architectures that decompose development into specialized roles (planning, implementation, testing, and review) offer a principled approach to iterative refinement, yet their empirical benefits remain underexplored.

We conduct a controlled comparison of four LLM-based code generation architectures on the HumanEval benchmark (164 tasks): a single-agent baseline, a multi-agent pipeline using one model across all roles, a multi-agent system with tiered models (1.5B--32B parameters) and difficulty-based routing, and an always-large variant. We further evaluate prompt repetition, a technique that duplicates user queries to improve token attention.

Multi-agent orchestration achieves 72.6\% pass rate versus 67.7\% for single-agent (+4.9 pp), with the review loop recovering 22\% of initially failing tasks. However, adaptive routing underperforms the always-large strategy (62.2\% vs 81.7\%), as the smallest model (1.5B) consistently fails and triggers costly escalation cascades. Prompt repetition degrades multi-agent accuracy by up to 7.9 pp, likely due to context dilution in already-rich conversation histories.

These results demonstrate that multi-agent pipelines improve correctness when paired with sufficiently capable models, but adaptive routing provides no benefit when tier-S models fall below a reliability threshold. The findings inform the design of cost-effective LLM orchestration systems.

\end{abstract}

